\documentclass[11pt,paper=a4,answers]{exam}
\usepackage{graphicx,lastpage}
\usepackage{upgreek}
\usepackage{censor}
\usepackage{gensymb}
\usepackage{amsmath}
\usepackage{multicol}
\usepackage{enumitem}
\usepackage{tasks}
\usepackage{adjustbox}
\usepackage{xcolor}
\usepackage{tfrupee}
\setlength{\voffset}{0.25in}
\usepackage{tabularx}
\newcolumntype{Y}{>{\centering\arraybackslash}X}
%==============================================================
% Customise Non-Essential Packages Below
% =============================================================
\usepackage[most]{tcolorbox}
\usepackage{eso-pic}
\usepackage{tikz}
\AddToShipoutPictureBG{%
\begin{tikzpicture}[remember picture,overlay]
\foreach \x in {-8,-4,0,4,8,12,16,20}
\foreach \y in {-8,-4,0,4,8,12,16,20} {
\node[opacity=0.05,rotate=45,anchor=center] at ([xshift=\x cm,yshift=\y cm]current page.center) {%
\begin{minipage}{2cm}
\centering
\includegraphics[width=2cm]{images/logo_black.png}\\
\small preppeo.com
\end{minipage}%
};
}
\end{tikzpicture}%
}
\printanswers

%==============================================================
% Customise Non-Essential Packages Above
% =============================================================
\definecolor{svgray}{rgb}{0.90, 0.90, 0.90}
\graphicspath{{images/}}

\usepackage[normalem]{ulem}
\renewcommand\ULthickness{1pt}   %%---> For changing thickness of underline
\setlength\ULdepth{3.5ex}%\maxdimen ---> For changing depth of underline
\renewcommand{\baselinestretch}{1}
\chead{
\includegraphics[scale=0.1,height=2cm ]{logo.png}
}
\pagestyle{empty}

\pagestyle{headandfoot}
\headrule
\cfoot{W: https://preppeo.com; M: +91-8130711689}

\begin{document}
%==============================================================
% Customise Question paper details below this
% =============================================================
\begin{center}
    \centering
    \renewcommand{\arraystretch}{1.5}
    \begin{tabularx}{\textwidth}{YYY}
    & \normalsize\bf{{\Large{{Quant}}}} & \\
    By Preppeo & \normalsize\bf{{alg\_sat\_med\_spr\_006}} & SAT\\
    \normalsize{{\textbf{{Algebra}} }} & \normalsize{{\textbf{{Linear equations, systems, inequalities }} }} & \normalsize{{\textbf{{Solving Linear Equations}} }} \\
    \end{tabularx}
\end{center}
\par\noindent\rule{\textwidth}{1pt}
% =============================================================
% Customise Question paper details above this
% =============================================================

\begin{questions}
\pointsinrightmargin
\pointsdroppedatright
\marksnotpoints
\pointpoints{mark}{marks}
\pointformat{\boldmath\themarginpoints}
\bracketedpoints

% =============================================================
% Question Begin
% =============================================================

\section*{SAT Algebra - Practice Set 4 (Medium)}

% Q1: Based on Rectangle Perimeter logic - WITH TIKZ
\question
The length of a rectangular garden is $60$ feet and the width is $w$ feet. The total perimeter of the garden is at most $280$ feet.
\begin{center}
\begin{tikzpicture}[scale=0.8]
\draw[thick, fill=green!10] (0,0) rectangle (4,2.5);
\node[below] at (2,0) {$60$};
\node[right] at (4,1.25) {$w$};
\node at (2,1.25) {\small Perimeter $\leq 280$};
\end{tikzpicture}
\end{center}
Which inequality represents this situation?
\begin{tasks}(2)
\task $2w + 60 \leq 280$
\task $2w + 120 \leq 280$
\task $2w + 120 \geq 280$
\task $w + 60 \leq 280$
\end{tasks}

\begin{solution}
The perimeter of a rectangle is the sum of all four sides, calculated as $P = 2(\text{length}) + 2(\text{width})$.
\vspace{5pt}
\begin{center}
Substitute the given values: $P = 2(60) + 2(w) = 120 + 2w$.
\end{center}
\vspace{5pt}
The phrase "at most $280$" indicates that the perimeter must be less than or equal to $280$.
\vspace{5pt}
Therefore, the correct inequality is $2w + 120 \leq 280$.
\vspace{5pt}
Answer: \colorbox{yellow!100}{B}
\end{solution}

\vspace{1cm}

% Q2: Based on Event Planner Budget logic
\question
A catering company charges a flat setup fee of $\$45$ plus $\$12.50$ per guest. The host has a maximum budget of $\$320$. What is the greatest number of guests possible without exceeding the budget?
\begin{tasks}(2)
\task $20$
\task $22$
\task $25$
\task $26$
\end{tasks}

\begin{solution}
Let $g$ represent the number of guests. We set up an inequality where the total cost is less than or equal to the budget.
\vspace{5pt}
\begin{center}
$45 + 12.50g \leq 320$
\vspace{5pt}
$12.50g \leq 320 - 45$
\vspace{5pt}
$12.50g \leq 275$
\vspace{5pt}
$g \leq \dfrac{275}{12.5} = 22$
\end{center}
\vspace{5pt}
The maximum number of guests the host can afford is $22$.
\vspace{5pt}
Answer: \colorbox{yellow!100}{B}
\end{solution}

\vspace{1cm}

% Q3: Based on Perpendicular Line logic - WITH TIKZ
\question
Line $m$ in the $xy$-plane is perpendicular to the line with the equation $y = 4$.
\begin{center}
\begin{tikzpicture}[scale=0.6]
\draw[thick,->] (-2,0) -- (3,0) node[right] {\tiny $x$};
\draw[thick,->] (0,-1) -- (0,4) node[above] {\tiny $y$};
\draw[blue, thick] (-2,3) -- (3,3) node[right] {\tiny $y=4$};
\draw[red, ultra thick] (1,-1) -- (1,4) node[above] {\tiny Line $m$};
\draw (1,3) rectangle (0.7,2.7);
\end{tikzpicture}
\end{center}
What is the slope of line $m$?
\begin{tasks}(2)
\task $0$
\task $-4$
\task $\frac{1}{4}$
\task The slope is undefined.
\end{tasks}

\begin{solution}
The line $y = 4$ is a horizontal line, which has a slope of $0$.
\vspace{5pt}
\begin{center}
Lines perpendicular to horizontal lines are vertical lines.
\end{center}
\vspace{5pt}
By definition, a vertical line has no horizontal change (run), which makes the denominator of the slope formula zero.
\vspace{5pt}
Therefore, the slope of a vertical line is undefined.
\vspace{5pt}
Answer: \colorbox{yellow!100}{D}
\end{solution}

\vspace{1cm}

% Q4: Based on Calories/Exercise logic
\question
In a fitness guide, it is estimated that a person burns $300$ calories for every $40$ minutes of swimming and $200$ calories for every $40$ minutes of jogging. A person has already completed $1$ hour of jogging. How many hours should they swim to burn a total of $1,200$ calories from both activities?
\begin{tasks}(2)
\task $2$
\task $3.5$
\task $4$
\task $4.5$
\task $1.5$
\end{tasks}

\begin{solution}
First, calculate the calories burned during the $1$ hour ($60$ minutes) of jogging.
\vspace{5pt}
\begin{center}
Jogging rate $= \dfrac{200 \text{ cal}}{40 \text{ min}} = 5 \text{ cal/min}$.
\vspace{5pt}
Calories from jogging $= 5 \times 60 = 300 \text{ calories}$.
\vspace{5pt}
Remaining calories needed $= 1,200 - 300 = 900 \text{ calories}$.
\end{center}
\vspace{5pt}
Next, find the swimming time ($t$) needed to burn $900$ calories.
\vspace{5pt}
\begin{center}
Swimming rate $= \dfrac{300 \text{ cal}}{40 \text{ min}} = 7.5 \text{ cal/min}$.
\vspace{5pt}
$7.5t = 900 \implies t = 120 \text{ minutes}$.
\vspace{5pt}
$120 \text{ minutes} = 2 \text{ hours}$.
\end{center}
\vspace{5pt}
Answer: \colorbox{yellow!100}{A}
\end{solution}

\vspace{1cm}

% Q5: Based on Infinitely Many Solutions logic - WITH TIKZ
\question
$-10x + 20y = 40$ \\
One of the two equations in a system is given above. The system has infinitely many solutions. Which of the following could be the second equation?
\begin{center}
\begin{tikzpicture}[scale=0.5]
\draw[step=1cm,gray!10,very thin] (-4,-2) grid (4,4);
\draw[thick,->] (-4,0) -- (4,0) node[right] {\tiny $x$};
\draw[thick,->] (0,-2) -- (0,4) node[above] {\tiny $y$};
\draw[ultra thick, blue] (-4,0) -- (4,4);
\node[blue] at (2,1) {\tiny Single Line for both eqns};
\end{tikzpicture}
\end{center}
\begin{tasks}(2)
\task $5x - 10y = -20$
\task $5x - 10y = 20$
\task $-5x - 10y = 20$
\task $10x + 20y = 40$
\end{tasks}

\begin{solution}
For a system to have infinitely many solutions, the two equations must represent the same line. This means one equation is a multiple of the other.
\vspace{5pt}
\begin{center}
Start with the given equation: $-10x + 20y = 40$.
\vspace{5pt}
Divide every term by $-2$:
\vspace{5pt}
$\dfrac{-10x}{-2} + \dfrac{20y}{-2} = \dfrac{40}{-2} \implies 5x - 10y = -20$.
\end{center}
\vspace{5pt}
This matches the equation in option A.
\vspace{5pt}
Answer: \colorbox{yellow!100}{A}
\end{solution}

\vspace{1cm}

% Q6: Based on value of x system logic
\question
$6x + y = 14$ \\
$y = 8x - 7$ \\
The solution to the given system of equations is $(x, y)$. What is the value of $x$?
\begin{tasks}(2)
\task $0.5$
\task $1.5$
\task $2$
\task $3.5$
\end{tasks}

\begin{solution}
We can solve for $x$ by substituting the expression for $y$ from the second equation into the first equation.
\vspace{5pt}
\begin{center}
$6x + (8x - 7) = 14$
\vspace{5pt}
$14x - 7 = 14$
\vspace{5pt}
$14x = 14 + 7 = 21$
\vspace{5pt}
$x = \dfrac{21}{14} = 1.5$
\end{center}
\vspace{5pt}
Answer: \colorbox{yellow!100}{B}
\end{solution}
% =============================================================
% Question End
% =============================================================

\end{questions}
\begin{center}
\rule{.5\textwidth}{1pt}
\end{center}
\end{document}
